\documentclass[12pt,a4paper]{article}

% ─── Packages ────────────────────────────────────────────────────────────────
\usepackage[top=2.5cm,bottom=2.5cm,left=3cm,right=3cm]{geometry}
\usepackage{graphicx}
\usepackage{xcolor}
\usepackage{listings}
\usepackage{longtable}
\usepackage{booktabs}
\usepackage{array}
\usepackage{tabularx}
\usepackage{hyperref}
\usepackage{enumitem}
\usepackage{fancyhdr}
\usepackage{titlesec}
\usepackage{tcolorbox}
\usepackage{fontenc}
\usepackage{inputenc}
\usepackage{lmodern}
\usepackage{parskip}
\usepackage{caption}
\usepackage{subcaption}
\usepackage{amsmath}
\usepackage{mdframed}

\tcbuselibrary{skins,breakable}

% ─── Colours ─────────────────────────────────────────────────────────────────
\definecolor{nanblue}{RGB}{30,64,175}
\definecolor{nangreen}{RGB}{22,101,52}
\definecolor{nanorange}{RGB}{154,52,18}
\definecolor{nangray}{RGB}{75,85,99}
\definecolor{codebg}{RGB}{243,244,246}
\definecolor{codefg}{RGB}{17,24,39}
\definecolor{borderblue}{RGB}{147,197,253}
\definecolor{lightblue}{RGB}{239,246,255}
\definecolor{lightyellow}{RGB}{254,252,232}
\definecolor{lightred}{RGB}{254,242,242}
\definecolor{lightgreen}{RGB}{240,253,244}

% ─── Hyperref ────────────────────────────────────────────────────────────────
\hypersetup{
    colorlinks=true,
    linkcolor=nanblue,
    urlcolor=nanblue,
    citecolor=nanblue,
    pdfauthor={Team NAN},
    pdftitle={NAN -- Computer Networks Phase 1 Report},
    pdfsubject={Computer Networks},
}

% ─── Header / Footer ─────────────────────────────────────────────────────────
\pagestyle{fancy}
\fancyhf{}
\fancyhead[L]{\small\textcolor{nangray}{NAN --- Computer Networks Phase 1}}
\fancyhead[R]{\small\textcolor{nangray}{\thepage}}
\fancyfoot[C]{\small\textcolor{nangray}{Team NAN $\cdot$ Lakshya $\cdot$ Kushal $\cdot$ Chinmay $\cdot$ Prachee}}
\renewcommand{\headrulewidth}{0.4pt}
\renewcommand{\footrulewidth}{0.4pt}

% ─── Section Formatting ──────────────────────────────────────────────────────
\titleformat{\section}
  {\Large\bfseries\color{nanblue}}
  {\thesection}{1em}{}[\textcolor{nanblue}{\rule{\linewidth}{0.5pt}}]

\titleformat{\subsection}
  {\large\bfseries\color{nangray}}
  {\thesubsection}{1em}{}

\titleformat{\subsubsection}
  {\normalsize\bfseries\color{nangray}}
  {\thesubsubsection}{1em}{}

% ─── Code Listing Style ───────────────────────────────────────────────────────
\lstset{
    backgroundcolor=\color{codebg},
    basicstyle=\ttfamily\footnotesize\color{codefg},
    breaklines=true,
    breakatwhitespace=true,
    frame=single,
    framesep=4pt,
    rulecolor=\color{borderblue},
    xleftmargin=8pt,
    xrightmargin=8pt,
    aboveskip=8pt,
    belowskip=8pt,
    showstringspaces=false,
    tabsize=4,
    keywordstyle=\color{nanblue}\bfseries,
    commentstyle=\color{nangreen}\itshape,
    stringstyle=\color{nanorange},
}

% ─── Custom Boxes ─────────────────────────────────────────────────────────────
\newtcolorbox{infobox}[1][]{
  colback=lightblue, colframe=nanblue,
  fonttitle=\bfseries\color{nanblue},
  title={#1}, breakable
}

\newtcolorbox{successbox}[1][]{
  colback=lightgreen, colframe=nangreen,
  fonttitle=\bfseries\color{nangreen},
  title={#1}, breakable
}

\newtcolorbox{warningbox}[1][]{
  colback=lightyellow, colframe=nanorange,
  fonttitle=\bfseries\color{nanorange},
  title={#1}, breakable
}

% ─── Title Page ──────────────────────────────────────────────────────────────
\begin{document}

\begin{titlepage}
  \centering
  \vspace*{1.5cm}

  {\Huge\bfseries\color{nanblue} NAN}\\[0.4cm]
  {\Large\bfseries\color{nangray} Private Network Service Platform}\\[0.2cm]
  {\large Computer Networks --- Phase 1 Report}

  \vspace{1cm}
  \textcolor{nanblue}{\rule{0.6\linewidth}{1.5pt}}
  \vspace{1cm}

  \begin{center}
    \begin{tabular}{ll}
      \toprule
      \textbf{Role} & \textbf{Team Member} \\
      \midrule
      DNS Server (Mac 1) & Lakshya \\
      Nginx Gateway (Mac 2) & Kushal \\
      Backend A (Mac 3) & Chinmay \\
      Backend B (Mac 4) & Prachee \\
      \bottomrule
    \end{tabular}
  \end{center}

  \vspace{1.2cm}
  {\large\textbf{Course:} Computer Networks}\\[0.2cm]
  {\large\textbf{Team:} NAN}\\[0.2cm]
  {\large\textbf{Phase:} 1 --- Build \& Observe}\\[0.2cm]
  {\large\textbf{Date:} October 2026}

  \vfill

  \begin{infobox}[Core Principle]
    \textit{``The application stays simple --- the network is the project.''}\\[0.3em]
    A client types a private domain name. It resolves through a team-owned DNS server,
    connects over HTTPS through a reverse proxy, and receives a response from one of
    two backend servers. Every packet is captured and explained.
  \end{infobox}

\end{titlepage}

% ─── Table of Contents ───────────────────────────────────────────────────────
\tableofcontents
\newpage

% ─────────────────────────────────────────────────────────────────────────────
\section{Project Overview}
% ─────────────────────────────────────────────────────────────────────────────

\subsection{Objective}

The NAN project builds a fully local, four-machine distributed service platform running on a private Wi-Fi LAN --- no cloud, no pre-configured servers. The goal is to demonstrate, measure, and explain every layer of a real network request: from DNS resolution, through TCP and TLS handshake, to HTTP load balancing and application response.

\subsection{Team and Machine Roles}

Each physical machine takes a dedicated network role. The table below maps each team member to their machine's responsibility and the software they operate.

\begin{table}[h!]
\centering
\caption{Team Machine Roles and IP Addresses}
\begin{tabular}{llllll}
\toprule
\textbf{Machine} & \textbf{Role} & \textbf{Member} & \textbf{IP Address} & \textbf{Port(s)} & \textbf{Software} \\
\midrule
Mac 1 & DNS Server & Lakshya & \texttt{10.7.3.114} & 53 & dnsmasq \\
Mac 2 & Nginx Gateway & Kushal & \texttt{10.7.29.219} & 443, 8080 & nginx + TLS \\
Mac 3 & Backend A & Chinmay & \texttt{10.7.20.70} & 3001 & Python HTTP \\
Mac 4 & Backend B & Prachee & \texttt{10.7.0.210} & 3002 & Python HTTP \\
\bottomrule
\end{tabular}
\end{table}

\subsection{DNS Records}

Two private domain names were configured, both resolving to the Nginx gateway:

\begin{lstlisting}
app.NAN.test  ->  10.7.29.219   (Nginx Gateway)
api.NAN.test  ->  10.7.29.219   (Nginx Gateway)
\end{lstlisting}

% ─────────────────────────────────────────────────────────────────────────────
\section{Network Architecture}
% ─────────────────────────────────────────────────────────────────────────────

\subsection{Topology}

All four machines connect to the same private Wi-Fi access point. There is no cloud component --- the entire network runs locally on the team's laptops.

\begin{figure}[h!]
\centering
\begin{mdframed}[backgroundcolor=codebg,linecolor=borderblue]
\begin{verbatim}
              Wi-Fi / LAN (private, ~10.7.0.0/16)
                           |
    +------------+---------+-----------+-----------+
    |            |                     |           |
    v            v                     v           v
[Router]   [DNS Server]         [Nginx Gateway] [Backend A] [Backend B]
           Lakshya              Kushal           Chinmay     Prachee
           10.7.3.114           10.7.29.219      10.7.20.70  10.7.0.210
           Port 53              Port 443/8080    Port 3001   Port 3002
           dnsmasq              nginx + TLS      Python      Python
\end{verbatim}
\end{mdframed}
\caption{Physical network topology}
\end{figure}

\subsection{Cloud Architecture Equivalents}

\begin{table}[h!]
\centering
\caption{Local components mapped to cloud equivalents}
\begin{tabular}{lll}
\toprule
\textbf{Local Component} & \textbf{AWS Equivalent} & \textbf{GCP Equivalent} \\
\midrule
dnsmasq on Mac 1 & Route 53 Private Hosted Zone & Cloud DNS \\
Nginx on Mac 2 & Application Load Balancer & HTTP(S) Load Balancer \\
Backend A on Mac 3 & EC2 Instance & Compute Engine VM \\
Backend B on Mac 4 & EC2 Instance & Compute Engine VM \\
Private LAN & VPC & VPC Network \\
Self-signed CA & AWS Certificate Manager & GCP Certificate Manager \\
\bottomrule
\end{tabular}
\end{table}

% ─────────────────────────────────────────────────────────────────────────────
\section{Task A --- LAN Setup}
% ─────────────────────────────────────────────────────────────────────────────

\subsection{Overview}

All four team machines were connected to the same private Wi-Fi network. Connectivity was verified by direct HTTP access from the gateway machine (Kushal, Mac 2) to both backend machines (Chinmay, Prachee).

\subsection{Evidence --- Direct Backend Access}

\begin{lstlisting}
=== Backend A ===
HTTP/1.0 200 OK
Server: BaseHTTP/0.6 Python/3.14.2
Content-Type: text/plain
Content-Length: 21
Cache-Control: max-age=60
X-Backend: A

Backend A is healthy

=== Backend B ===
HTTP/1.0 200 OK
Server: BaseHTTP/0.6 Python/3.14.7
Content-Type: text/plain
Content-Length: 21
Cache-Control: max-age=60
X-Backend: B

Backend B is healthy
\end{lstlisting}

\begin{successbox}[Result]
Direct HTTP connectivity from Mac 2 to both Mac 3 (\texttt{10.7.20.70:3001}) and Mac 4 (\texttt{10.7.0.210:3002}) confirmed. Both backends returned \texttt{200 OK} with the correct \texttt{X-Backend} identifier.
\end{successbox}

% ─────────────────────────────────────────────────────────────────────────────
\section{Task B --- Private DNS Server}
% ─────────────────────────────────────────────────────────────────────────────

\subsection{Configuration}

dnsmasq was installed on Lakshya's machine (Mac 1, \texttt{10.7.3.114}) and configured to resolve the \texttt{.NAN.test} zone while forwarding all other queries to public resolvers (8.8.8.8, 1.1.1.1), preserving normal internet connectivity.

\begin{lstlisting}[language=bash, caption=dns/dnsmasq.conf]
port=53
listen-address=10.7.3.114
bind-interfaces

# NAN project DNS records
address=/app.NAN.test/10.7.29.219
address=/api.NAN.test/10.7.29.219

# Forward external DNS to public resolvers
server=8.8.8.8
server=1.1.1.1

log-queries
log-facility=-
\end{lstlisting}

\subsection{Evidence --- DNS Resolution}

\begin{lstlisting}[caption=dig output for app.NAN.test]
;; Got answer:
;; ->>HEADER<<- opcode: QUERY, status: NOERROR, id: 6248

;; QUESTION SECTION:
;app.NAN.test.      IN  A

;; ANSWER SECTION:
app.NAN.test.  0   IN  A   10.7.29.219

;; SERVER: 10.7.3.114#53(10.7.3.114)
;; MSG SIZE rcvd: 57
\end{lstlisting}

\begin{lstlisting}[caption=dig output for api.NAN.test]
;; ANSWER SECTION:
api.NAN.test.  0   IN  A   10.7.29.219

;; SERVER: 10.7.3.114#53(10.7.3.114)
\end{lstlisting}

\begin{successbox}[Result]
Both \texttt{app.NAN.test} and \texttt{api.NAN.test} resolve to \texttt{10.7.29.219} (the Nginx gateway) using the team's private DNS server at \texttt{10.7.3.114:53}. Status: \texttt{NOERROR}. Response is authoritative (\texttt{aa} flag set).
\end{successbox}

\subsection{DNS vs.\ IP Layer Distinction}

A key conceptual point demonstrated: DNS resolution and TCP/IP connectivity are independent layers.

\begin{table}[h!]
\centering
\caption{Layer independence demonstrated through DNS failure scenarios}
\begin{tabular}{lll}
\toprule
\textbf{Condition} & \textbf{DNS Result} & \textbf{TCP/IP Result} \\
\midrule
Correct DNS configured & Resolves correctly & Connection succeeds \\
Wrong DNS configured & Fails to resolve & Direct IP still reachable \\
DNS record wrong IP & Resolves (wrong IP) & Connection to wrong host \\
DNS server offline & Name lookup fails & Direct IP still reachable \\
\bottomrule
\end{tabular}
\end{table}

% ─────────────────────────────────────────────────────────────────────────────
\section{Task C --- Backend Services}
% ─────────────────────────────────────────────────────────────────────────────

\subsection{Design}

Two minimal Python HTTP servers were built. Both bind to \texttt{0.0.0.0} so they are reachable from all machines on the LAN. Each exposes two endpoints and identifies itself in every response via the \texttt{X-Backend} header.

\begin{table}[h!]
\centering
\caption{Backend API endpoints}
\begin{tabular}{lll}
\toprule
\textbf{Endpoint} & \textbf{Response Body} & \textbf{Headers} \\
\midrule
\texttt{GET /} & \texttt{Backend A} / \texttt{Backend B} & \texttt{X-Backend: A} or \texttt{B} \\
\texttt{GET /api/status} & \texttt{Backend A is healthy} & \texttt{X-Backend: A}, \texttt{Cache-Control: max-age=60} \\
\bottomrule
\end{tabular}
\end{table}

\subsection{Source Code --- Backend A}

\begin{lstlisting}[language=Python, caption=backends/backend-a/server.py]
from http.server import BaseHTTPRequestHandler, HTTPServer

HOST = "0.0.0.0"
PORT = 3001
BACKEND = "A"

class Handler(BaseHTTPRequestHandler):

    def do_GET(self):
        if self.path == "/api/status":
            body = f"Backend {BACKEND} is healthy\n".encode()
            self.send_response(200)
            self.send_header("Content-Type", "text/plain")
            self.send_header("Content-Length", str(len(body)))
            self.send_header("Cache-Control", "max-age=60")
            self.send_header("X-Backend", BACKEND)
            self.end_headers()
            self.wfile.write(body)
        else:
            body = f"Backend {BACKEND}\n".encode()
            self.send_response(200)
            self.send_header("Content-Type", "text/plain")
            self.send_header("Content-Length", str(len(body)))
            self.send_header("X-Backend", BACKEND)
            self.end_headers()
            self.wfile.write(body)

server = HTTPServer((HOST, PORT), Handler)
print(f"Backend {BACKEND} listening on {HOST}:{PORT}")
server.serve_forever()
\end{lstlisting}

Backend B is identical with \texttt{BACKEND = "B"} and \texttt{PORT = 3002}.

% ─────────────────────────────────────────────────────────────────────────────
\section{Task D --- Nginx Reverse Proxy and Load Balancer}
% ─────────────────────────────────────────────────────────────────────────────

\subsection{Configuration}

Nginx on Mac 2 (\texttt{10.7.29.219}) acts as the single entry point. Clients never connect directly to backend IPs. Round-robin load balancing distributes requests across both backends.

\begin{lstlisting}[caption=nginx/nan.conf (abridged)]
upstream nan_backends {
    server 10.7.20.70:3001;   # Backend A - Chinmay
    server 10.7.0.210:3002;   # Backend B - Prachee
}

server {
    listen 443 ssl;
    server_name app.NAN.test api.NAN.test;

    ssl_certificate     /opt/homebrew/etc/nginx/nan-server.crt;
    ssl_certificate_key /opt/homebrew/etc/nginx/nan-server.key;

    ssl_protocols       TLSv1.2 TLSv1.3;
    ssl_ciphers         HIGH:!aNULL:!MD5;

    location / {
        proxy_pass http://nan_backends;
        proxy_set_header Host              $host;
        proxy_set_header X-Real-IP         $remote_addr;
        proxy_set_header X-Forwarded-Proto $scheme;
        proxy_next_upstream error timeout http_502;
    }
}
\end{lstlisting}

\subsection{Traffic Path}

\begin{figure}[h!]
\centering
\begin{mdframed}[backgroundcolor=codebg,linecolor=borderblue]
\begin{verbatim}
Client
  |
  |  HTTPS  app.NAN.test:443
  v
Nginx Gateway  10.7.29.219:443
  |
  |  Round-Robin  (plain HTTP on LAN)
  |
  +----[odd]----->  Backend A  10.7.20.70:3001
  |
  +----[even]---->  Backend B  10.7.0.210:3002
\end{verbatim}
\end{mdframed}
\caption{Request routing through Nginx to backends}
\end{figure}

\subsection{Evidence --- Load Balancing}

Twelve consecutive HTTPS requests to \texttt{app.NAN.test} produced the following pattern:

\begin{lstlisting}
Request 1:  X-Backend: B        Request  7:  X-Backend: B
Request 2:  X-Backend: A        Request  8:  X-Backend: A
Request 3:  X-Backend: B        Request  9:  X-Backend: B
Request 4:  X-Backend: A        Request 10:  X-Backend: A
Request 5:  X-Backend: B        Request 11:  X-Backend: B
Request 6:  X-Backend: A        Request 12:  X-Backend: A
\end{lstlisting}

\begin{successbox}[Result]
Perfect round-robin distribution: Backend B served odd requests, Backend A served even requests, alternating across all 12 requests without interruption.
\end{successbox}

% ─────────────────────────────────────────────────────────────────────────────
\section{Task E --- HTTPS and TLS}
% ─────────────────────────────────────────────────────────────────────────────

\subsection{Certificate Authority Setup}

A local private Certificate Authority named \textbf{NAN Local CA} was created using OpenSSL. The CA certificate was distributed to all client machines and added to the macOS system trust store, enabling certificate verification without the \texttt{-k} flag.

\begin{lstlisting}[language=bash, caption=Adding NAN Local CA to macOS trust store]
sudo security add-trusted-cert -d -r trustRoot \
    -k /Library/Keychains/System.keychain nan-ca.crt
\end{lstlisting}

\subsection{Server Certificate}

\begin{table}[h!]
\centering
\caption{Server certificate properties}
\begin{tabular}{ll}
\toprule
\textbf{Property} & \textbf{Value} \\
\midrule
Subject (CN) & \texttt{app.NAN.test} \\
Issuer & \texttt{NAN Local CA} \\
Organisation & \texttt{NAN} \\
Country / State / City & IN / West Bengal / Kolkata \\
Valid From & Oct 4, 2026 \\
Valid Until & Oct 4, 2027 \\
TLS Version & TLSv1.3 \\
Cipher Suite & \texttt{AEAD-CHACHA20-POLY1305-SHA256} \\
\bottomrule
\end{tabular}
\end{table}

\subsection{TLS Termination Architecture}

TLS is terminated at Nginx. Backend communication uses plain HTTP on the trusted private LAN.

\begin{figure}[h!]
\centering
\begin{mdframed}[backgroundcolor=codebg,linecolor=borderblue]
\begin{verbatim}
Client <------ TLS 1.3 encrypted ------> Nginx (10.7.29.219:443)
                                               |
                                      plain HTTP on LAN
                                               |
                                 +-------------+-------------+
                           Backend A                   Backend B
                        10.7.20.70:3001             10.7.0.210:3002
\end{verbatim}
\end{mdframed}
\caption{TLS termination at Nginx edge}
\end{figure}

\subsection{TLS Handshake Observed}

The \texttt{curl -v} output captured the complete TLS 1.3 handshake:

\begin{lstlisting}
* (OUT), TLS handshake, Client hello (1)
* (IN),  TLS handshake, Server hello (2)
* (IN),  TLS handshake, Unknown (8)       <- TLS 1.3 extensions
* (IN),  TLS handshake, Certificate (11)
* (IN),  TLS handshake, CERT verify (15)
* (IN),  TLS handshake, Finished (20)
* (OUT), TLS handshake, Finished (20)
* SSL connection using TLSv1.3 / AEAD-CHACHA20-POLY1305-SHA256
* SSL certificate verify ok.
\end{lstlisting}

\begin{table}[h!]
\centering
\caption{TLS 1.3 handshake message sequence}
\begin{tabular}{lll}
\toprule
\textbf{Message} & \textbf{Direction} & \textbf{Purpose} \\
\midrule
ClientHello & Client $\rightarrow$ Server & Advertises TLS versions, ciphers, SNI \\
ServerHello & Server $\rightarrow$ Client & Selects TLS 1.3, CHACHA20 cipher \\
Certificate & Server $\rightarrow$ Client & Sends \texttt{nan-server.crt} \\
CertVerify & Server $\rightarrow$ Client & Proves possession of private key \\
Finished & Both directions & Confirms handshake; session keys active \\
\bottomrule
\end{tabular}
\end{table}

After \textit{Finished}, all application data is encrypted. HTTP payload is not visible in Wireshark --- only opaque TLS Application Data records appear.

\begin{successbox}[Result]
HTTPS established without certificate warnings. \texttt{SSL certificate verify ok.} confirmed by curl. TLS 1.3 with CHACHA20-POLY1305 cipher in use.
\end{successbox}

% ─────────────────────────────────────────────────────────────────────────────
\section{Task F --- HTTP Caching}
% ─────────────────────────────────────────────────────────────────────────────

\subsection{Implementation}

The \texttt{/api/status} endpoint includes a \texttt{Cache-Control: max-age=60} header in every response. This instructs clients and intermediate caches to treat the response as fresh for 60 seconds.

\subsection{Evidence}

\begin{lstlisting}[caption=curl -sI https://api.NAN.test/api/status]
HTTP/1.1 200 OK
Server: nginx/1.31.6
Content-Type: text/plain
Content-Length: 21
Connection: keep-alive
Cache-Control: max-age=60
X-Backend: B
\end{lstlisting}

\subsection{Cache Behaviour Explained}

\begin{table}[h!]
\centering
\caption{HTTP cache states}
\begin{tabular}{lp{5.5cm}l}
\toprule
\textbf{State} & \textbf{Behaviour} & \textbf{HTTP Status} \\
\midrule
Fresh hit & Response served from local cache; no network request made & (none) \\
Conditional request & Client re-validates with \texttt{If-None-Match} or \texttt{If-Modified-Since} & \texttt{304 Not Modified} \\
Full new request & Cache expired or flushed; complete round-trip to backend & \texttt{200 OK} \\
\bottomrule
\end{tabular}
\end{table}

% ─────────────────────────────────────────────────────────────────────────────
\section{Task G --- Complete Protocol Flow Capture}
% ─────────────────────────────────────────────────────────────────────────────

\subsection{OSI Layer Mapping}

Every layer of the OSI model is exercised in a single client request:

\begin{table}[h!]
\centering
\caption{OSI layer mapping for the NAN request}
\begin{tabular}{cllp{5cm}}
\toprule
\textbf{Layer} & \textbf{Name} & \textbf{Protocol} & \textbf{Where / What} \\
\midrule
7 & Application & DNS, HTTP/1.1 & dnsmasq, Python backends, nginx \\
6 & Presentation & TLS 1.3 & Nginx: encrypts/decrypts \\
5 & Session & TLS session & Key exchange, cipher negotiation \\
4 & Transport & TCP (443, 3001, 3002) / UDP (53) & All machines \\
3 & Network & IPv4 (\texttt{10.7.x.x}) & Wi-Fi LAN routing \\
2 & Data Link & 802.11 Wi-Fi frames & NIC on each machine \\
1 & Physical & Wi-Fi radio signal & Access point \\
\bottomrule
\end{tabular}
\end{table}

\subsection{Port Identification}

\begin{table}[h!]
\centering
\caption{Ports and socket pairs in a single request}
\begin{tabular}{lllll}
\toprule
\textbf{Traffic} & \textbf{Protocol} & \textbf{Src Port} & \textbf{Dst Port} & \textbf{Machine} \\
\midrule
DNS query & UDP & Ephemeral & 53 & Client $\rightarrow$ Mac 1 \\
DNS response & UDP & 53 & Ephemeral & Mac 1 $\rightarrow$ Client \\
HTTPS & TCP & Ephemeral & 443 & Client $\rightarrow$ Mac 2 \\
Proxy to Backend A & TCP & Ephemeral & 3001 & Mac 2 $\rightarrow$ Mac 3 \\
Proxy to Backend B & TCP & Ephemeral & 3002 & Mac 2 $\rightarrow$ Mac 4 \\
\bottomrule
\end{tabular}
\end{table}

\subsection{Wireshark Captures}

\begin{table}[h!]
\centering
\caption{Wireshark capture files}
\begin{tabular}{ll}
\toprule
\textbf{File} & \textbf{Contents} \\
\midrule
\texttt{NAN-phase1-dns-capture.pcapng} & DNS query/response for \texttt{app.NAN.test} \\
\texttt{NAN\_phase1\_baseline\_traffic.pcapng} & Full system: DNS + TCP + TLS + HTTP \\
\texttt{NAN\_phase1\_backend\_failure\_failover.pcapng} & Backend A down, failover to B \\
\bottomrule
\end{tabular}
\end{table}

\textbf{Key Wireshark filters used during analysis:}

\begin{lstlisting}[language=bash]
dns                         # DNS query and response packets
tcp.flags.syn == 1          # TCP three-way handshake
tls.handshake               # TLS ClientHello, ServerHello, Certificate
tcp.port == 443             # All HTTPS traffic
tcp.flags.reset == 1        # TCP RST (connection refused)
ip.addr == 10.7.3.114       # DNS server traffic
\end{lstlisting}

\subsection{TCP Three-Way Handshake}

Observed in the Wireshark baseline capture (\texttt{tcp.flags.syn == 1 \&\& tcp.port == 443}):

\begin{lstlisting}
Client        -->  [SYN]      seq=X             --> 10.7.29.219:443
10.7.29.219   -->  [SYN-ACK]  seq=Y, ack=X+1   --> Client
Client        -->  [ACK]      ack=Y+1            --> 10.7.29.219:443
\end{lstlisting}

The handshake completes before any TLS or application data is transmitted, confirming TCP's reliable, connection-oriented nature.

% ─────────────────────────────────────────────────────────────────────────────
\section{Phase 1 Failure Demonstrations}
% ─────────────────────────────────────────────────────────────────────────────

\begin{table}[h!]
\centering
\caption{Phase 1 mandatory failure scenarios}
\begin{tabular}{p{4cm}p{4.5cm}p{4cm}}
\toprule
\textbf{Failure Injected} & \textbf{Observed Behaviour} & \textbf{Layer Affected} \\
\midrule
Wrong DNS resolver on client & Name resolution fails; direct IP still reachable & Application (DNS, L7) \\
DNS record points to wrong IP & DNS succeeds; TCP connection reaches wrong host & Network (L3) \\
One backend stopped & Nginx routes all traffic to remaining backend & Application upstream \\
Both backends stopped & Nginx returns \texttt{502 Bad Gateway} & Application upstream \\
Wrong destination port & TCP \texttt{Connection refused}; host reachable & Transport (TCP, L4) \\
\bottomrule
\end{tabular}
\end{table}

\subsection{Backend Failover --- Detailed}

When Backend A (Mac 3) was stopped:

\begin{lstlisting}
Before failure:  X-Backend: A, B, A, B, A, B ...
After Backend A stopped:  X-Backend: B, B, B, B, B ...
After Backend A restarted:  X-Backend: A, B, A, B ...
\end{lstlisting}

Nginx's passive health check (\texttt{proxy\_next\_upstream error timeout}) detected the TCP connection failure to \texttt{10.7.20.70:3001} and automatically routed subsequent requests to \texttt{10.7.0.210:3002} (Backend B). The client received no errors during the failover.

\subsection{502 Bad Gateway}

When both backends were stopped simultaneously, Nginx returned:

\begin{lstlisting}
HTTP/1.1 502 Bad Gateway
Server: nginx/1.31.6
\end{lstlisting}

This confirms the boundary between the edge proxy layer and the application backend layer: the TLS handshake still succeeds (Nginx is reachable), but the upstream HTTP connection cannot be established.

% ─────────────────────────────────────────────────────────────────────────────
\section{Security Notes}
% ─────────────────────────────────────────────────────────────────────────────

\begin{warningbox}[Private Key Handling]
The following files contain private cryptographic material and are \textbf{not committed} to the repository. They reside only on the gateway machine:
\begin{itemize}
  \item \texttt{nan-server.key} --- Nginx TLS private key
  \item \texttt{nan-ca.key} --- NAN Local CA private key
\end{itemize}
A \texttt{.gitignore} rule (\texttt{*.key}, \texttt{*.pem}) prevents accidental commit.
\end{warningbox}

Only the public server certificate (\texttt{tls/nan-server.crt}) and public CA certificate are distributed and committed.

% ─────────────────────────────────────────────────────────────────────────────
\section{Project File Structure}
% ─────────────────────────────────────────────────────────────────────────────

\begin{lstlisting}
cn-phase1/
|-- README.md
|-- architecture/
|   |-- topology.md          Network diagram, IP table, cloud equivalents
|   `-- request-flow.md      Step-by-step protocol flow, OSI mapping
|-- backends/
|   |-- backend-a/
|   |   `-- server.py        Backend A (port 3001)
|   `-- backend-b/
|       `-- server.py        Backend B (port 3002)
|-- dns/
|   `-- dnsmasq.conf         dnsmasq config with .NAN.test records
|-- nginx/
|   `-- nan.conf             Reverse proxy, load balancer, TLS config
|-- tls/
|   |-- README.md            TLS setup, CA details, trust store
|   `-- nan-server.crt       Public server certificate
|-- scripts/
|   |-- verify.sh            Pre-demo automated checks
|   `-- failure-demo.sh      Failure scenario commands
|-- docs/
|   `-- submission-notes.md  Task completion status, checklist
`-- evidence/
    |-- 01-lan/              LAN connectivity proof
    |-- 02-dns/              DNS resolution output
    |-- 03-https/            TLS handshake verification
    |-- 04-load-balancing/   Round-robin load balancing log
    |-- 05-cache/            Cache-Control header demonstration
    |-- 06-wireshark/        Wireshark captures and screenshots
    `-- 07-failure/          Backend failure and failover evidence
\end{lstlisting}

% ─────────────────────────────────────────────────────────────────────────────
\section{Summary}
% ─────────────────────────────────────────────────────────────────────────────

\begin{table}[h!]
\centering
\caption{Phase 1 task completion summary}
\begin{tabular}{clll}
\toprule
\textbf{Task} & \textbf{Description} & \textbf{Status} & \textbf{Evidence} \\
\midrule
A & Establish Private LAN & \textcolor{nangreen}{\textbf{Complete}} & \texttt{evidence/01-lan/} \\
B & Configure Private DNS & \textcolor{nangreen}{\textbf{Complete}} & \texttt{evidence/02-dns/} \\
C & Build Two Backend Services & \textcolor{nangreen}{\textbf{Complete}} & \texttt{backends/} \\
D & Nginx Reverse Proxy + Load Balancer & \textcolor{nangreen}{\textbf{Complete}} & \texttt{evidence/04-load-balancing/} \\
E & HTTPS / TLS & \textcolor{nangreen}{\textbf{Complete}} & \texttt{evidence/03-https/} \\
F & HTTP Caching & \textcolor{nangreen}{\textbf{Complete}} & \texttt{evidence/05-cache/} \\
G & Complete Protocol Flow Capture & \textcolor{nangreen}{\textbf{Complete}} & \texttt{evidence/06-wireshark/} \\
\midrule
\multicolumn{2}{l}{Failure Scenarios} & \textcolor{nangreen}{\textbf{Complete}} & \texttt{evidence/07-failure/} \\
\bottomrule
\end{tabular}
\end{table}

All Phase 1 mandatory tasks are complete. The NAN platform successfully demonstrates the full DNS $\rightarrow$ TCP $\rightarrow$ TLS $\rightarrow$ HTTP $\rightarrow$ Load Balancer $\rightarrow$ Backend stack on a private local network, with Wireshark evidence at every layer and controlled failure scenarios showing layer independence.

% ─────────────────────────────────────────────────────────────────────────────
\vfill
\begin{center}
\textcolor{nangray}{\rule{0.5\linewidth}{0.4pt}}\\[0.5em]
{\small \textcolor{nangray}{Team NAN $\cdot$ Computer Networks Phase 1 $\cdot$ October 2026}}\\
{\small \textcolor{nangray}{Lakshya $\cdot$ Kushal $\cdot$ Chinmay $\cdot$ Prachee}}
\end{center}

\end{document}
